% Revised Introduction addressing reviewer comments:
% each requested topic is discussed in a separate paragraph.
% Test list and chamber features are deferred to Section 3 (Problem Definition).

\section{Introduction}
\label{sec:intro}

The refrigerator industry operates in a highly competitive global market characterized by compressed innovation cycles and increasingly stringent market-specific requirements for performance, safety, and energy efficiency \citep{IEA2021,IEC60335-1,IEC60335-2-24,IEC62552}. Since products cannot proceed to certification, production, or commercialization until design approval tests have been successfully completed, research \& development (R\&D) testing has become a critical operational bottleneck. Delays in testing postpone market entry and reduce firms' responsiveness to changing market conditions, while inefficient utilization of capital-intensive testing infrastructure diminishes the return on R\&D investment \citep{OECD_MSTI2024,OECD_MSTI2025}. Furthermore, launch timing and quality assurance are closely linked, as the trade-off between accelerated time-to-market and product quality has a direct impact on commercial performance \citep{KimKoenigsbergOfek2022}. This study is motivated by refrigerator R\&D testing operations at a white goods manufacturer located in Turkey.

Refrigerator R\&D testing comprises a multi-stage portfolio of laboratory tests performed under controlled environmental conditions. The catalog includes gas amount determination, pull-down, energy-efficiency measurement, performance evaluation at climatic-class temperatures, optional freezing-capacity and temperature-rise tests, and consumer-usage tests. Because refrigerators must demonstrate compliance under the ambient conditions of their target markets, tests are executed at different temperature set-points (commonly spanning \(10^\circ\mathrm{C}\) to \(43^\circ\mathrm{C}\)), under ambient or humidity-controlled conditions (including \(85\%\) relative humidity for selected consumer-usage tests), and, when required by regional electrical standards, under voltage-compatible chamber settings. These changing conditions exist because climatic class, energy regulation, and electrical specification differ across regions; a single ambient environment cannot certify a product for all markets. Individual tests are long relative to daily planning cycles: typical durations range from about 5 days (pull-down and temperature rise) to 10--15 days (gas determination, performance, energy efficiency, and consumer usage). Consequently, each assignment occupies scarce chamber capacity for an extended period, and sequencing mistakes cannot be corrected quickly. The detailed test list and durations are presented in Section~\ref{sec:pd}.

Which tests are required for a given product, and when they must be completed, are determined by commercial and regulatory inputs rather than by a fixed universal checklist. Target-market region and climatic class select the applicable performance and energy-efficiency variants; electrical specifications determine whether voltage-compatible execution is mandatory; product engineers may add optional tests; and project due dates reflect planned launch windows and certification milestones. Scheduling must therefore coordinate heterogeneous, product-specific test mixes under shared due-date pressure. Sample-size determination is an integral part of this decision process. In current practice, a uniform three-sample policy is often used because it is simple to administer. With few samples, upstream pull-down workload remains limited, but downstream tests that could otherwise run in parallel on different physical prototypes are serialized, which can increase completion times for due-date-critical projects. With several samples, downstream parallelism improves and tardiness may decrease, yet each additional sample multiplies pull-down operations and intensifies congestion in temperature-constrained environments. The net effect is therefore non-monotone: increasing samples does not guarantee better on-time performance, and the value of an extra sample depends on due dates, required test mixes, chamber eligibility, and stage-based release logic.

The supporting testing infrastructure consists of heterogeneous climate chambers, each containing multiple stations that operate under a shared chamber-level environment. Chambers differ in station capacity and in their technical ability to adjust temperature, humidity, and voltage. Temperature control is generally available, whereas humidity and voltage adjustment are restricted to selected chambers; thus, not every test can be executed in every chamber. Because all stations in a chamber share one environment, chamber configuration and station-level scheduling are tightly coupled. Although chambers can be reconfigured, frequent environment changes are undesirable: temperature and humidity transitions require stabilization time, interrupt ongoing tests that need a stable ambient setting, and reduce effective capacity. In practice, chamber environments are therefore treated as planning decisions that should remain stable over the scheduling horizon rather than as freely changing machine states. Chamber capabilities and the chamber--station structure are detailed in Section~\ref{sec:pd}.

The operational problem we seek to solve is the following. Given a portfolio of refrigerator development projects with product-specific test requirements and due dates, decide (i)~how many physical samples to allocate to each project, (ii)~which environmental configuration to assign to each chamber, and (iii)~how to schedule all resulting test operations on chamber stations so that environmental, humidity, voltage, station-capacity, sample-availability, and stage-based precedence constraints are respected. The primary goal is to minimize total tardiness; among solutions with equal total tardiness, lower total sample usage is preferred. We refer to this integrated planning problem as the \emph{Test Chamber Scheduling with Sample Optimization Problem} (TCSSOP). Unlike classical scheduling problems with a fixed job set, sample decisions endogenously generate part of the workload and simultaneously create or restrict parallelization opportunities.

TCSSOP is challenging because its main decisions are interdependent and the induced workload is decision-dependent. Sample sizes change both the volume of pull-down work and the feasible degree of downstream parallelism; chamber-environment assignments determine which workloads can be processed where; and stage-based precedence means that early congestion propagates into later stages. Heterogeneous chamber capabilities create eligibility bottlenecks for humidity- and voltage-restricted tests, while long test durations amplify the cost of poor allocations. From an operations research perspective, the setting combines resource-constrained multi-project scheduling, unrelated parallel resources with eligibility restrictions, and tardiness objectives---features that are already computationally demanding and typically require specialized heuristic methods \citep{HartmannBriskorn2022,GomezSanchezLallaRuizFernandezGilCastroVoss2023,UlagaDurasevicJakobovic2022,MaeckerShenMonch2023,AfzaliradRezaeian2016}. Existing integrated planning streams do not fully cover this structure: order acceptance and scheduling coordinate selection with sequencing but generally treat processing workload as exogenous \citep{Slotnick2011,WangYe2019}, whereas reliability demonstration testing optimizes sample size under statistical criteria without detailed chamber-level scheduling \citep{Fernandez2022}.

This study makes four main contributions. First, it benchmarks fixed uniform sample policies against endogenous sample-size optimization and quantifies their effect on total tardiness. Second, it includes an equal-budget control experiment in which a uniform six-sample policy uses the same average prototype consumption as the optimized solution, thereby isolating the effect of project-level allocation from total sample quantity. Third, it develops an iterative framework that coordinates sample-size decisions, chamber-environment assignment, and feasible scheduling. Fourth, it introduces a Variable Neighborhood Search (VNS) diversification strategy that perturbs workload-defining sample-size decisions rather than only sequencing or assignment decisions.

To solve TCSSOP, we propose an iterative framework. For a given sample vector, chamber environments are first generated through workload-balancing logic and further refined by local search. Based on these resulting chamber settings, sample sizes are improved through bounded local search, with each candidate evaluated by an Earliest Due Date (EDD) based constructive scheduler that explicitly validates feasibility. Since sample-size decisions change the workload distribution across environments, chamber assignments are re-optimized iteratively. When local improvement stagnates, a VNS diversification mechanism perturbs sample sizes and re-optimizes the resulting workload structure.

The remainder of the paper is organized as follows. We review the related studies in Section~\ref{sec:literature_review}. We define the problem in Section~\ref{sec:pd} and present the notation used throughout the paper, including the test catalog and chamber features. We present the solution methodology proposed to address the problem in Section~\ref{sec:solution}. We report the results of the computational experiments in Section~\ref{sec:computational_study}. We conclude the paper with our findings in Section~\ref{sec:conc}.
